\chapter{FUNDAMENTAÇÃO TEÓRICA}

A base teórica deste trabalho divide-se em dois grandes eixos. O primeiro trata da produção do espaço e da crítica à economia política, tendo como principal referência a obra de \cite{dask1}. O autor demonstra como o capitalismo produz e reproduz suas próprias contradições no espaço urbano e nas relações sociais.

O segundo eixo aborda as práticas educativas e o papel da avaliação no processo de ensino-aprendizagem. \cite{luckesi-avaprend} contribui para essa discussão ao evidenciar que a avaliação escolar, frequentemente, tem sido utilizada como um instrumento de exclusão e classificação, em vez de ser um mecanismo de diagnóstico e transformação da prática pedagógica.